% !TeX program = xelatex
\documentclass[12pt,a4paper]{article}
\usepackage{fontspec}
\setmainfont{Times New Roman}
\usepackage[english,russian]{babel}
\usepackage{amsmath,amssymb,mathrsfs,cancel}
\usepackage[margin=22mm]{geometry}
\usepackage{graphicx,float}
\usepackage[unicode,hidelinks]{hyperref}
\DeclareMathOperator{\grad}{grad}
\DeclareMathOperator{\Div}{div}
\newcommand{\Rey}{\mathrm{Re}}
\setlength{\emergencystretch}{3em}
\sloppy
\interfootnotelinepenalty=10000
\newcommand{\topic}[1]{\par\medskip\noindent\textbf{#1}\par\nobreak\smallskip}
\renewcommand{\thefootnote}{\arabic{footnote})}
\begin{document}
\section*{Вопрос 10}
Описание многофазной фильтрации. Обобщенный закон Дарси, относительная фазовая проницаемость. Учет капиллярной разности давлений между фазами, функция Леверетта.

\medskip
\noindent\textbf{Источник:} К. С. Басниев, А. М. Власов, И. Н. Кочина, В. М. Максимов. \textit{Подземная гидравлика}. Москва: Недра, 1986. Глава 1, § 5, с. 19--22; глава 9, §§ 2--3, с. 199--204.
\par\bigskip
\begin{quote}\sloppy
Углеводородные системы могут быть гомо- и гетерогенными. В гомогенной системе все ее части имеют одинаковые физические и химические свойства. Составляющие гомогенной системы (называемые компонентами) «размазаны» по всему пространству и взаимодействуют на молекулярном уровне. Для гетерогенной системы физические и химические свойства в разных точках различны. Гетерогенные системы состоят из фаз. Фаза --- это часть системы, которая является гомогенной и отделена от других фаз отчетливыми границами. Смесь воды, нефти и газа в пласте --- типичный пример гетерогенной среды.

Главными характеристиками движения многофазной системы являются насыщенность и скорость фильтрации каждой фазы.

Насыщенностью $\sigma_i$ порового пространства $i$-й фазой называется доля объема пор $\Delta V_i$, занятая этой фазой в элементарном объеме:
\begin{equation}\sigma_i=\Delta V_i/\Delta V_\text{п},\quad i=1,2,\ldots,n.\tag{9.1}\label{q10:eq-9-1}\end{equation}
Очевидно, что
\begin{equation}\sum_{i=1}^n\sigma_i=1.\tag{9.2}\label{q10:eq-9-2}\end{equation}
Таким образом, в $n$-фазной системе имеется $(n-1)$ независимая насыщенность. В частности, при исследовании фильтрации смеси двух фаз используется лишь одна из насыщенностей, которая обозначается в дальнейшем $\sigma$ (обычно это насыщенность вытесняющей фазы).

Движение каждой из фаз характеризуется вектором скорости фильтрации $\vec w_i$ данной фазы, который (по аналогии со скоростью фильтрации однородной жидкости) определяется как вектор, проекция которого на некоторое направление $l$ равна отношению объемного расхода $Q_i$ данной фазы к площадке $\Omega_l$, перпендикулярной к указанному направлению:
\begin{equation}(\vec w_i)_l=Q_i/\Omega_l,\quad i=1,2,\ldots,n.\tag{9.3}\label{q10:eq-9-3}\end{equation}
Площадка $\Omega_l$ пересекает как твердую, так и подвижные фазы.
\par\nobreak\smallskip{\small К. С. Басниев, А. М. Власов, И. Н. Кочина, В. М. Максимов. \textit{Подземная гидравлика}. Москва: Недра, 1986. Глава 9, § 2, с. 199--200.}\par
\end{quote}

\begin{quote}\sloppy
Экспериментально установлено, что расход каждой фазы растет с увеличением перепада давления и насыщенности данной фазой, а закон фильтрации каждой из фаз по аналогии с законом Дарси можно записать в виде
\begin{equation}w_1=\frac{k_1^*\Delta p_1}{\mu_1L},\qquad w_2=\frac{k_2^*\Delta p_2}{\mu_2L},\tag{1.18}\label{q10:eq-1-18}\end{equation}
где $w_1$ и $w_2$ --- скорости фильтрации фаз; $\mu_1$ и $\mu_2$ --- динамические коэффициенты вязкости жидкостей; $\Delta p_1$ и $\Delta p_2$ --- разности давлений в соответствующих фазах; $k_1^*$ и $k_2^*$ --- фазовые проницаемости. Величины $k_i^*$ ($i=1,2$) являются для $i$-й жидкости проницаемостями в обычном смысле в условиях совместной фильтрации. Эти величины зависят от природы пористой среды (и, прежде всего, от ее абсолютной проницаемости $k$, определяемой по данным о фильтрации однородной жидкости), а также от насыщенности пористой среды каждой фазой. При описании двухфазных течений обычно вместо фазовых проницаемостей $k_i^*$ вводят так называемые
«относительные проницаемости» $k_i$ фаз, определяемые из отношений
\begin{equation}k_1=k_1^*/k,\qquad k_2=k_2^*/k.\tag{1.19}\label{q10:eq-1-19}\end{equation}
В большинстве опытов показано, что для данной структуры пористой среды относительные проницаемости $k_i$ являются в основном функциями насыщенности, а если и наблюдается влияние иных параметров (например, отношения коэффициентов вязкости $\mu_0$ фаз), то им обычно пренебрегают. Тогда с учетом (1.19) закон Дарси (1.18) для каждой из фаз записывается в виде
\begin{equation}w_i=k\frac{k_i(\sigma)\Delta p_i}{\mu_iL},\qquad i=1,2.\tag{1.20}\label{q10:eq-1-20}\end{equation}
\par\nobreak\smallskip{\small К. С. Басниев, А. М. Власов, И. Н. Кочина, В. М. Максимов. \textit{Подземная гидравлика}. Москва: Недра, 1986. Глава 1, § 5, с. 19--20.}\par
\end{quote}

\begin{quote}\sloppy
Типичный вид экспериментальных кривых фазовых проницаемостей $k_i(\sigma)$ приведен на рис. 1.6 (кривая \textit{1} относится к более смачивающей жидкости --- воде, кривая \textit{2} --- к менее смачивающей --- нефти; кривая $1'$ относится к случаю, когда первая фаза является газом и $\sigma$ --- газонасыщенность). Отметим некоторые характерные особенности этих кривых. Для каждой фазы существует предельная насыщенность ($\sigma_*$ и $1-\sigma^*$) такая, что при меньших значениях насыщенности эта фаза неподвижна. Движение первой фазы может происходить только в том случае, если $\sigma>\sigma_*$ (для водонефтяной системы $\sigma_*$ называют насыщенностью связанной водой). Для второй фазы связанная насыщенность равна $1-\sigma^*$ и называется остаточной нефтенасыщенностью. Таким образом, совместное течение двух фаз имеет место лишь в следующем интервале изменения насыщенности жидкостью 1:
\[\sigma_*\le\sigma\le\sigma^*.\]
Сумма относительных проницаемостей для каждого фиксированного значения $\sigma$ меньше 1:
\[k_1(\sigma)+k_2(\sigma)<1,\qquad0<\sigma<1.\]
\par\nobreak\smallskip{\small К. С. Басниев, А. М. Власов, И. Н. Кочина, В. М. Максимов. \textit{Подземная гидравлика}. Москва: Недра, 1986. Глава 1, § 5, с. 20.}\par
\end{quote}

\begin{quote}\sloppy
При совместном течении двух фаз в пористой среде, по крайней мере, одна из них образует систему, граничащую со скелетом породы и частично с другой жидкостью. Из-за избирательного смачивания твердой породы одной из жидкостей площадь контакта каждой из фаз со скелетом пористой среды значительно превышает площадь контакта фаз между собой. Это позволяет предположить, что каждая фаза движется по занятым ею поровым каналам под действием «своего» давления независимо от других фаз, т. е. так, как если бы она была ограничена только твердыми стенками. При этом, естественно, сопротивление, испытываемое каждой фазой при совместном течении, отлично от того, которое было бы при фильтрации только одной из них. Опыты показывают, что расход каждой фазы растет с увеличением насыщенности и градиента давления. Закон фильтрации каждой из фаз при учете силы тяжести по аналогии с законом Дарси можно записать в следующем виде:
\begin{equation}\vec w_i=-\frac{k}{\mu_i}k_i(\sigma)(\grad p_i-\rho_i\vec g).\tag{9.4}\label{q10:eq-9-4}\end{equation}
Здесь $k$ --- абсолютная проницаемость пласта, определяемая по данным о фильтрации однородной жидкости; $\mu_i$ --- коэффициент динамической вязкости фаз; $p_i$ --- давление в фазах; $\rho_i$ --- плотность фаз; $\vec g$ --- вектор ускорения свободного падения; $k_i(\sigma)$ --- относительные фазовые проницаемости, определяемые экспериментально; $\sigma$ --- насыщенность одной из фаз.

Будем считать для определенности, что $\sigma=\sigma_1$ --- насыщенность вытесняющей (или более смачивающей) фазы. Тогда из (9.2) имеем $\sigma_2=1-\sigma$. Понятие относительной фазовой проницаемости $k_i(\sigma)$, введенное в гл. 1, играет важную роль при изучении совместного течения нескольких жидкостей в пористой среде. Мы будем исходить из условия, что относительные проницаемости являются однозначными функциями насыщенностей (см. рис. 1.6).

Характерные особенности многофазной фильтрации связаны также с влиянием поверхностного натяжения. Давления в фазах $p_1$ и $p_2$, вообще говоря, не равны друг другу из-за капиллярных эффектов, приводящих к скачку давления на границе раздела фаз, так что
\begin{equation}p_2-p_1=p_\text{к},\tag{9.5}\label{q10:eq-9-5}\end{equation}
где $p_\text{к}$ --- капиллярное давление (или капиллярный скачок).

Большее давление будет на стороне жидкости, не смачивающей твердые зерна породы.

Будем предполагать, как это принято, что капиллярное давление при совместном течении жидкостей совпадает с капиллярным давлением в равновесном состоянии для того же значения насыщенности и при одном и том же направлении ее изменения (увеличении или уменьшении). Поэтому будем считать, что при движении капиллярное давление можно представить в виде известной экспериментальной функции насыщенности (рис. 9.1):
\begin{equation}p_\text{к}=p_\text{к}(\sigma)=\alpha_\text{п}\cos\theta\sqrt{\frac{m}{k}}J(\sigma),\tag{9.6}\label{q10:eq-9-6}\end{equation}
где $\alpha_\text{п}$ --- коэффициент межфазного поверхностного натяжения; $\theta$ --- статический краевой угол смачивания между жидкостями и породой; $m$ --- пористость; $J(\sigma)$ --- безразмерная функция Леверетта.


\par\nobreak\smallskip{\small К. С. Басниев, А. М. Власов, И. Н. Кочина, В. М. Максимов. \textit{Подземная гидравлика}. Москва: Недра, 1986. Глава 9, § 2, с. 200--201.}\par
\end{quote}

\begin{quote}\sloppy
Будем для простоты рассматривать совместное изотермическое течение двух фаз в однородной пористой среде без фазовых переходов и химических реакций. Система уравнений, описывающая совместную фильтрацию фаз, строится на основе уравнений неразрывности для каждой фазы, уравнений движения (закона фильтрации) и соответствующих замыкающих соотношений.
\par\nobreak\smallskip{\small К. С. Басниев, А. М. Власов, И. Н. Кочина, В. М. Максимов. \textit{Подземная гидравлика}. Москва: Недра, 1986. Глава 9, § 3, с. 202.}\par
\end{quote}

\begin{quote}\sloppy
или в векторной форме
\begin{equation}\frac{\partial}{\partial t}(m\rho_1\sigma)+\Div(\rho_1\vec w_1)=0.\tag{9.9}\label{q10:eq-9-9}\end{equation}
Аналогично выводится уравнение сохранения массы для второй фазы:
\begin{equation}\frac{\partial}{\partial t}[m\rho_2(1-\sigma)]+\Div(\rho_2\vec w_2)=0.\tag{9.10}\label{q10:eq-9-10}\end{equation}
\par\nobreak\smallskip{\small К. С. Басниев, А. М. Власов, И. Н. Кочина, В. М. Максимов. \textit{Подземная гидравлика}. Москва: Недра, 1986. Глава 9, § 3, с. 203.}\par
\end{quote}

\begin{quote}\sloppy
Если жидкости и пористую среду можно предполагать несжимаемыми, то вместо уравнений (9.9) и (9.10) имеем
\begin{equation}m\frac{\partial\sigma}{\partial t}+\Div\vec w_1=0,\qquad-m\frac{\partial\sigma}{\partial t}+\Div\vec w_2=0.\tag{9.11}\label{q10:eq-9-11}\end{equation}
\topic{Уравнения движения для многофазной фильтрации}
При записи закона фильтрации будем предполагать, что в любой точке каждая из фаз находится в термодинамически равновесном состоянии. Тогда для течения двухфазной смеси можно ввести в рассмотрение относительные проницаемости $k_i(\sigma)$ и капиллярное давление $\rho_\text{к}(\sigma)$, зависящее только от насыщенности.
Будем рассматривать только однонаправленные процессы фильтрации, не учитывая гистерезисных явлений.

Тогда выполняется закон фильтрации (9.4):
\begin{equation}\vec w_i=-\frac{k}{\mu_i}k_i(\sigma)(\grad p_i-\rho_i\vec g),\tag{9.12}\label{q10:eq-9-12}\end{equation}
а связь между давлениями в фазах определяется равенствами (9.5) и (9.6):
\begin{equation}p_2-p_1=p_\text{к}(\sigma)=\alpha_\text{п}\cos\theta\sqrt{\frac{m}{k}}J(\sigma).\tag{9.13}\label{q10:eq-9-13}\end{equation}
Для замыкания полученной системы уравнений необходимо задать связи параметров, характеризующих свойства фаз и пористой среды, с давлением.

\topic{Уравнения состояния флюидов}
При изотермических условиях фильтрации плотность и вязкость каждой из фаз определяются давлением в данной фазе:
\begin{equation}\rho_i=\rho_i(p_i),\qquad\mu_i=\mu_i(p_i),\qquad(i=1,2).\tag{9.14}\label{q10:eq-9-14}\end{equation}
Пласт будем считать недеформируемым.

Таким образом, получена замкнутая система уравнений (9.9), (9.10), (9.12), (9.13) и (9.14) для определения всех неизвестных параметров: насыщенности $\sigma$, давления $p_i$, скорости фильтрации $\vec w_i$, плотности $\rho_i$ и вязкости $\mu_i$ фаз.
\par\nobreak\smallskip{\small К. С. Басниев, А. М. Власов, И. Н. Кочина, В. М. Максимов. \textit{Подземная гидравлика}. Москва: Недра, 1986. Глава 9, § 3, с. 204.}\par
\end{quote}

\end{document}
